\documentclass[twoside]{article}

\usepackage[preprint]{aistats2027}
\usepackage[round,authoryear]{natbib}
\usepackage{amsmath,amssymb,amsthm}
\usepackage{mathtools}
\usepackage{microtype}
\usepackage{booktabs}
\usepackage{algorithm}
\usepackage{algorithmic}
\usepackage{url}

\renewcommand{\bibname}{References}
\renewcommand{\bibsection}{\subsubsection*{\bibname}}
\bibliographystyle{apalike}

\newtheorem{theorem}{Theorem}
\newtheorem{proposition}[theorem]{Proposition}
\newtheorem{corollary}[theorem]{Corollary}
\newtheorem{lemma}[theorem]{Lemma}
\newtheorem{remark}[theorem]{Remark}
\newcommand{\E}{\mathbb E}
\newcommand{\Prb}{\mathbb P}
\newcommand{\one}{\mathbf 1}
\newcommand{\TV}{d_{\rm TV}}
\newcommand{\diag}{\operatorname{diag}}
\newcommand{\supp}{\operatorname{supp}}
\newcommand{\ip}[2]{\langle #1,#2\rangle}

\begin{document}
\runningtitle{Confidence Cost of Pairwise Learning}
\runningauthor{}

\twocolumn[
\aistatstitle{The Confidence Cost\\of Pairwise Distribution Learning}

\aistatsauthor{Pahan Dewasurendra}

\aistatsaddress{Johns Hopkins University\\[4pt] July 30, 2026}
]

\begin{abstract}
Pairwise conditional access to an unknown distribution returns a draw from
the distribution restricted to any requested pair. At constant confidence,
this weak oracle is as sample-efficient for labeled total-variation learning
as ordinary sampling, up to constants. We show that this equivalence fails at
high confidence.

For a distribution on $n$ labels, accuracy $\varepsilon$, and failure
probability $\delta$, we prove the exact minimax query complexity
\[
 \Theta\!\left(\frac{n\log(1/\delta)}{\varepsilon^2}\right).
\]
The corresponding ordinary-sampling complexity is
$\Theta((n+\log(1/\delta))/\varepsilon^2)$. Thus confidence amplification
costs up to a factor $n$ under pairwise access.

The upper bound reuses one trajectory of a pairwise-winner Markov chain whose
stationary law is the target. Although its worst-case relaxation time is
linear in $n$, we prove that its empirical distribution has the i.i.d.-scale
stationary risk $O(\sqrt{n/T})$. The proof compares its conductances with a
minimum-mass graph. After sorting the target probabilities, this graph has an
explicit Helmert eigensystem, and a capped-mass Green-function bound
telescopes without any minimum-mass or dynamic-range assumption. A geometric
median of independent trajectories gives the high-confidence result. The
lower bound uses a heavy label and uniform light labels. Every adaptive pair
query then carries only $O(\varepsilon^2/n)$ KL information, which makes the
multiplicative confidence cost unavoidable.
\end{abstract}

\section{Introduction}

Suppose an unknown distribution $p$ on $[n]$ cannot be sampled directly.
Instead, a learner chooses two labels $i,j$ and observes a draw from $p$
conditioned on $\{i,j\}$. This pairwise conditional, or PCOND, oracle is the
sampling form of the Bradley--Terry--Luce choice model
\citep{bradley1952rank,luce1959individual}. It appears when providers expose
only local menus, when comparisons are easier to elicit than calibrated
scores, and when global samples cannot be observed.

Pairwise access identifies every labeled distribution, including targets
with arbitrarily large dynamic range. The quantitative picture was recently
clarified by \citet{kleinberg2026providers}. Their structural theory of
restricted providers gives
\begin{equation}
 \Omega\!\left(\frac{n+\log(1/\delta)}{\varepsilon^2}\right)
 \ \le q_{\rm pair}\ \le
 O\!\left(\frac{n\log^2 n\log(n/\delta)}{\varepsilon^2}\right),
 \label{eq:prior-gap}
\end{equation}
and asks whether the remaining polylogarithmic gap can be closed. Existing
comparison estimators do not settle this question. Their useful uniform
bounds generally impose a bounded parameter range, a fixed passive design,
or a different loss on log qualities
\citep{shah2015estimation,agarwal2018accelerated,
hendrickx2019resistance,hendrickx2020minimax}.

We close the gap, but the answer differs from the ordinary-sampling
benchmark. Pairwise learning has a multiplicative, not additive, dependence
on confidence.

\begin{theorem}[Minimax pairwise learning]
\label{thm:minimax}
There are universal constants $c,C>0$ such that, for $n\ge4$,
$0<\varepsilon\le1/16$, and $0<\delta\le1/8$,
\begin{equation}
 c\frac{n\log(1/\delta)}{\varepsilon^2}
 \le q_{\rm pair}(n,\varepsilon,\delta)
 \le C\frac{n\log(1/\delta)}{\varepsilon^2}.
 \label{eq:minimax}
\end{equation}
The upper bound is attained by a polynomial-time adaptive learner and allows
arbitrary zero probabilities.
\end{theorem}

For comparison, ordinary i.i.d. sampling has complexity
$\Theta((n+\log(1/\delta))/\varepsilon^2)$
\citep{canonne2020short}. At constant confidence the two models both require
$\Theta(n/\varepsilon^2)$ observations. If
$\log(1/\delta)\gg n$, pairwise access requires a factor $\Theta(n)$ more.
Table~\ref{tab:rates} summarizes the separation.

\begin{table}[t]
\centering
\small
\caption{Minimax labeled TV learning rates.}
\label{tab:rates}
\setlength{\tabcolsep}{5pt}
\begin{tabular}{@{}lc@{}}
\toprule
Observation model & query or sample complexity \\
\midrule
Ordinary samples & $\Theta((n+\log(1/\delta))/\varepsilon^2)$ \\
Pairwise conditional & $\Theta(n\log(1/\delta)/\varepsilon^2)$ \\
Prior pairwise upper & $O(n\log^2n\log(n/\delta)/\varepsilon^2)$ \\
\bottomrule
\end{tabular}
\end{table}

\paragraph{Why mixing time is misleading.}
The upper bound starts from a natural chain. From state $i$, propose a
uniformly random $j\ne i$, query $\{i,j\}$, and retain the returned winner.
Its stationary law is $p$. Restarting after each $O(n)$-step burn-in gives a
quadratic learner. A single trajectory has worst-case relaxation time
$\Theta(n)$, so a black-box spectral-gap bound still suggests quadratic
cost. The slow mode, however, occurs only when its stationary mass is small
or concentrated in a few coordinates. Total variation automatically weights
these modes by their statistical relevance.

Our main technical result makes this compensation exact. We dominate the
chain's Green function by that of the complete graph with edge conductance
$\min(p_i,p_j)/n$. Once $p_1\ge\cdots\ge p_m>0$, this graph is diagonalized by
the Helmert contrasts. Its eigenvalues are capped masses
\[
 c_k/n,\qquad c_k=kp_k+\sum_{j>k}p_j
                  =\sum_j\min\{p_j,p_k\}.
\]
The coordinate Green functions then split into left, local, and right terms.
Each aggregate term is bounded by one telescoping inequality. This yields
$O(\sqrt{n/T})$ TV risk uniformly over the whole simplex.

\paragraph{Why confidence is different.}
Take one label of mass near $1/2$ and spread the remaining mass uniformly
over $n-1$ labels. Changing the heavy mass by $\Theta(\varepsilon)$ changes
the outcome probability of every informative pair by only
$\Theta(\varepsilon/n)$, while that outcome itself has probability
$\Theta(1/n)$. Each query therefore contributes only
$O(\varepsilon^2/n)$ KL divergence. This obstruction applies transcript by
transcript to adaptive learners and forces the product
$n\log(1/\delta)/\varepsilon^2$.

\paragraph{Contributions.}
Our results are entirely theoretical.
\begin{enumerate}
\item We determine the minimax PCOND complexity at every confidence level,
closing the dimension and accuracy gap in \eqref{eq:prior-gap}.
\item We prove a distribution-free occupation-risk theorem for the winner
chain. The result has no lower mass, bounded ratio, or warm-start assumption.
\item We identify a confidence separation between pairwise conditional and
ordinary sampling, and prove that adaptive pair selection cannot remove it.
\end{enumerate}

\section{Model and Algorithm}
\label{sec:model}

Let $\Delta_n$ be the simplex on $[n]$. For $p\in\Delta_n$, a query to an
unordered pair $\{i,j\}$ returns $i$ with probability
$p_i/(p_i+p_j)$ and $j$ otherwise. If both masses vanish, the response may
follow any fixed convention. A learner may choose each pair using its past
queries, responses, and private randomness.

We define $q_{\rm pair}(n,\varepsilon,\delta)$ as the least fixed query
budget for which some learner returns $\widehat p$ satisfying
\begin{align*}
 \sup_{p\in\Delta_n}
 \Prb_p\{\TV(\widehat p,p)>\varepsilon\}&\le\delta,\\
 \TV(p,q)&=\tfrac12\|p-q\|_1.
\end{align*}

One base estimator follows a single winner-chain trajectory. Repeating it
and aggregating by a metric median gives Algorithm~\ref{alg:learner}. An
$L_1$ geometric median is any minimizer over $\Delta_n$ of the sum of
distances to the candidates. It is computable by a linear program.

\begin{algorithm}[t]
\caption{High-confidence pairwise learner}
\label{alg:learner}
\begin{algorithmic}[1]
\STATE Set $C_0=2+2\sqrt2$,
$B=\lceil(n-1)\log(48/\varepsilon)\rceil$,
$T=\lceil(48^2C_0^2n/\varepsilon^2\rceil$, and
$R=\lceil18\log(1/\delta)\rceil$.
\FOR{$r=1,\ldots,R$}
  \STATE Start at an arbitrary $X_0\in[n]$.
  \FOR{$t=0,\ldots,B+T-1$}
    \STATE Draw $J_t$ uniformly from $[n]\setminus\{X_t\}$.
    \STATE Query $\{X_t,J_t\}$ and set $X_{t+1}$ to its winner.
  \ENDFOR
  \STATE Let $\widehat p_r$ be the empirical law of
  $X_{B+1},\ldots,X_{B+T}$.
\ENDFOR
\RETURN an $L_1$ geometric median of $\widehat p_1,\ldots,\widehat p_R$.
\end{algorithmic}
\end{algorithm}

The stated constants simplify the proof and are not optimized. Since
$\log(48/\varepsilon)=O(1/\varepsilon^2)$, the total query count has the
order in Theorem~\ref{thm:minimax}.

\subsection{What the algorithm changes}
\label{sec:algorithm-comparison}

The global chain itself is not new. It appears in the generic complete-family
learner of \citet{kleinberg2026providers}. The difference is which states are
kept and how they are analyzed. That learner runs the chain for
$B=\Theta(n\log(1/\varepsilon))$ steps, keeps only the terminal state, and
restarts independently $M=\Theta((n+\log(1/\delta))/\varepsilon^2)$ times.
Even in the pairwise case, this costs
\[
 \Theta\!\left(
 \frac{n(n+\log(1/\delta))\log(1/\varepsilon)}{\varepsilon^2}
 \right)
\]
queries. Their faster hierarchical learner keeps trajectories inside many
tree cells, then reconstructs $p$ from binary split estimates. It pays two
tree-depth logarithms and a simultaneous-confidence logarithm.

Algorithm~\ref{alg:learner} instead burns in one global chain and retains
every later state. This removes repeated mixing, but standard MCMC analysis
does not justify the improvement. The uniform minorization in
\eqref{eq:minorization} gives gap only $1/(n-1)$. Applying that gap to each
coordinate and summing gives the unhelpful bound $O(n/\sqrt T)$. The next
section replaces this worst-mode argument by a joint coordinate calculation.

\section{One Trajectory Has i.i.d.-Scale Risk}
\label{sec:risk}

We first analyze a stationary trajectory. On the positive support of $p$,
the winner chain has transition kernel
\begin{equation}
 P(i,j)=\frac1{n-1}\frac{p_j}{p_i+p_j},\quad i\ne j,
 \label{eq:kernel}
\end{equation}
with the diagonal chosen to make rows sum to one. Detailed balance holds
because
\begin{equation}
 p_iP(i,j)=\frac{p_ip_j}{(n-1)(p_i+p_j)}=p_jP(j,i).
 \label{eq:balance}
\end{equation}

\begin{theorem}[Stationary occupation risk]
\label{thm:risk}
Let $(X_t)_{t=1}^T$ be a stationary trajectory of \eqref{eq:kernel}, and let
$\widehat p_T$ be its empirical law. For every $p\in\Delta_n$,
\begin{equation}
 \E_p\TV(\widehat p_T,p)
 \le (2+2\sqrt2)\sqrt{\frac nT}.
 \label{eq:risk}
\end{equation}
\end{theorem}

The result is not a consequence of the worst-case gap, which can be
$\Theta(1/n)$. We prove it through coordinate-specific Green functions.
Let $f_i(x)=\one\{x=i\}-p_i$ and
\begin{equation}
 H_i=\ip{f_i}{(I-P)^{-1}f_i}_{p},
 \label{eq:green}
\end{equation}
where the inverse acts on mean-zero functions in $L_2(p)$.

There is a useful electrical interpretation. Let $D=\diag(p)$ and let
$L=D(I-P)$ be the weighted graph Laplacian associated with
\eqref{eq:balance}. Since
\[
 Df_i=p_i(e_i-p)=:b_i,
\]
the coordinate Green function is
\begin{equation}
 H_i=b_i^\top L^\dagger b_i.
 \label{eq:electrical}
\end{equation}
Thus $H_i$ is the minimum energy needed to send source
$p_i(1-p_i)$ from $i$ and absorb $p_ip_j$ at every $j\ne i$. A small global
gap can coexist with small coordinate energy because each source is scaled
by its target mass.

\begin{lemma}[Positive spectrum and variance]
\label{lem:variance}
The kernel $P$ is positive semidefinite on $L_2(p)$. Moreover,
\begin{equation}
 \operatorname{Var}_p\!\left(\frac1T\sum_{t=1}^Tf_i(X_t)\right)
 \le\frac{2H_i}{T}.
 \label{eq:variance}
\end{equation}
\end{lemma}

For positive semidefiniteness, use
\[
 \frac{p_ip_j}{p_i+p_j}(f_i-f_j)^2
 \le p_if_i^2+p_jf_j^2
\]
and sum over pairs. Reversibility then gives eigenvalues in $[0,1]$.
Expanding an additive functional in the eigenbasis and summing its geometric
covariances proves \eqref{eq:variance}.

It remains to control all $H_i$ jointly. Let $m=|\supp(p)|$ and relabel the
positive masses so that $p_1\ge\cdots\ge p_m>0$. Define
\begin{equation}
 c_k=kp_k+\sum_{j>k}p_j.
 \label{eq:capped}
\end{equation}

\begin{lemma}[Capped-mass Green bound]
\label{lem:profile}
For $i\in[m]$, define
\begin{align}
 A_i={}&\sum_{k=2}^{i-1}\frac{(1-c_k)^2}{k(k-1)c_k}
 +\one\{i\ge2\}\frac{(i-c_i)^2}{i(i-1)c_i}\notag\\
 &+\sum_{k=i+1}^{m}\frac{c_k}{k(k-1)}.
 \label{eq:Ai}
\end{align}
Then
\begin{equation}
 H_i\le2np_i^2A_i,
 \qquad
 \sum_{i=1}^{m}p_i\sqrt{A_i}\le2+2\sqrt2.
 \label{eq:profile}
\end{equation}
\end{lemma}

\paragraph{Proof idea.}
The Dirichlet conductance in \eqref{eq:balance} obeys
\begin{equation}
 \frac{p_ip_j}{(n-1)(p_i+p_j)}
 \ge\frac{\min\{p_i,p_j\}}{2n}.
 \label{eq:conductance-domination}
\end{equation}
Let $L_0$ be the graph Laplacian with conductance
$\min\{p_i,p_j\}/n$. The Helmert vectors
\begin{equation}
 u_k=\frac{(\underbrace{1,\ldots,1}_{k-1},-(k-1),0,\ldots,0)}
 {\sqrt{k(k-1)}}
 \label{eq:helmert}
\end{equation}
diagonalize $L_0$, with $L_0u_k=(c_k/n)u_k$. Projecting the coordinate source
$p_i(e_i-p)$ onto this basis gives \eqref{eq:Ai} exactly. The factor two in
\eqref{eq:profile} follows from \eqref{eq:conductance-domination}.

For completeness, the projections are explicit. One has
\[
 u_k^\top p=\frac{1-c_k}{\sqrt{k(k-1)}},
\]
and hence
\begin{equation}
 \frac{u_k^\top b_i}{p_i}=
 \begin{cases}
 -(1-c_k)/\sqrt{k(k-1)},&k<i,\\
 -(i-c_i)/\sqrt{i(i-1)},&k=i,\\
 c_k/\sqrt{k(k-1)},&k>i.
 \end{cases}
 \label{eq:projections}
\end{equation}
Dividing their squares by $c_k/n$ produces the three terms of
\eqref{eq:Ai}. Since the exact conductances dominate half of those of
$L_0$, the inverse-Laplacian variational principle yields
$L^\dagger\preceq2L_0^\dagger$ on the zero-sum subspace. This proves the
first part of Lemma~\ref{lem:profile}.

For the aggregate bound, split $A_i$ into the three sums in
\eqref{eq:Ai}. Jensen's inequality and a change in summation order bound the
first and third contributions by one each. For example,
\[
 \sum_i p_i\sqrt{\sum_{k<i}
 \frac{(1-c_k)^2}{k(k-1)c_k}}
 \le\sqrt{\sum_{k\ge2}\frac1{k(k-1)}}\le1,
\]
where $\sum_{i>k}p_i\le c_k$. For the middle term, let
$S_i=\sum_{j\ge i}p_j$. Since $c_i\ge S_i$,
\begin{align*}
 \sum_{i\ge2}p_i\frac{i-c_i}{\sqrt{i(i-1)c_i}}
 &\le\sqrt2\sum_i\frac{p_i}{\sqrt{S_i}}\\
 &\le2\sqrt2.
\end{align*}
The last inequality telescopes using
$(S_i-S_{i+1})/\sqrt{S_i}\le2(\sqrt{S_i}-\sqrt{S_{i+1}})$.
The supplement gives every projection and boundary case.

\begin{corollary}[Profile-aware risk]
\label{cor:profile-risk}
With $A_i$ as in \eqref{eq:Ai}, every stationary trajectory satisfies
\begin{align}
 \E\TV(\widehat p_T,p)
 &\le\sqrt{\frac nT}\,\Psi(p),\notag\\
 \Psi(p)&:=\sum_{i\in\supp(p)}p_i\sqrt{A_i}
 \le2+2\sqrt2.
 \label{eq:profile-risk}
\end{align}
\end{corollary}

\paragraph{Proof of Theorem~\ref{thm:risk}.}
By Lemmas~\ref{lem:variance} and \ref{lem:profile},
\begin{align*}
 \E\TV(\widehat p_T,p)
 &\le\frac12\sum_i
 \sqrt{\operatorname{Var}(\widehat p_T(i))}\\
 &\le\sqrt{\frac nT}\sum_i p_i\sqrt{A_i}
 \le(2+2\sqrt2)\sqrt{\frac nT}.
\end{align*}

\subsection{Examples and the slow mode}
\label{sec:examples}

The profile functional in \eqref{eq:profile-risk} reveals why a global gap
bound is loose. If $p$ is uniform on all $n$ labels, then $c_k=1$ and direct
substitution gives
\[
 A_i=1-\frac1n,\qquad
 \Psi(p)=\sqrt{1-\frac1n}.
\]
The winner chain is a lazy complete-graph walk. Its TV occupation risk is
therefore bounded by $\sqrt{(n-1)/T}$, at the ordinary empirical-distribution
scale.

At the other extreme, take
\begin{equation}
 p_1=\frac12,\qquad p_j=\frac1{2(n-1)},\quad j\ge2.
 \label{eq:half-heavy}
\end{equation}
The indicator $Y_t=\one\{X_t=1\}$ is itself a two-state Markov chain. From
either aggregate state it flips with probability exactly $1/n$. Its
nonconstant eigenvalue is $1-2/n$, so estimating the heavy mass from one
trajectory has variance of order $n/T$. This is a genuine slow mode.
Nevertheless, there is only one such aggregate coordinate. The remaining
$n-1$ light coordinates leave their individual states at constant rate. The
sum of all coordinate standard deviations remains $O(\sqrt{n/T})$, exactly
as Lemma~\ref{lem:profile} certifies.

For a nearly point-mass target, the slow aggregate mode becomes still slower,
but its stationary variance shrinks in the same proportion. In capped-mass
language, the local denominator $c_k$ decreases only where the corresponding
tail mass also decreases. The left and middle telescopes in the proof are the
finite-dimensional expression of this compensation.

These examples also separate estimation from independent sampling. The
uniform target behaves nearly as if the retained states were independent.
The half-heavy target has only about $T/n$ effective observations of its
heavy-versus-light split. That split is precisely the source of the
high-confidence lower bound in Section~\ref{sec:lower}.

\section{Burn-in and Confidence Amplification}
\label{sec:upper}

The winner chain admits the pointwise minorization
\begin{equation}
 P(i,\cdot)\ge\frac1{n-1}p(\cdot).
 \label{eq:minorization}
\end{equation}
This also holds from a zero-mass state. Hence, from every initial law $\mu$,
\begin{equation}
 \TV(\mu P^B,p)\le e^{-B/(n-1)}.
 \label{eq:burn}
\end{equation}
Applying the same Markov kernel to two initial states cannot increase TV, so
\eqref{eq:burn} also bounds the distance between the entire post-burn path
and a stationary path.

With the $B,T$ in Algorithm~\ref{alg:learner}, Theorem~\ref{thm:risk} and
\eqref{eq:burn} give
\[
 \E\TV(\widehat p_r,p)\le\varepsilon/24.
\]
Markov's inequality shows that each candidate is within $\varepsilon/4$ with
probability at least $5/6$. Hoeffding's inequality implies that at least
$2R/3$ candidates are good except with probability $e^{-R/18}\le\delta$.

If $G$ of $R$ points in any metric space lie within radius $r$ of $p$, where
$G>R/2$, every geometric median $q$ satisfies
\begin{equation}
 d(q,p)\le\frac{2G}{2G-R}r.
 \label{eq:median}
\end{equation}
Indeed, compare its sum of distances with the sum at $p$ and apply the
triangle inequality separately to the good and bad points. With
$G\ge2R/3$ and $r=\varepsilon/4$, \eqref{eq:median} proves the upper bound in
Theorem~\ref{thm:minimax}.

\paragraph{Computation.}
The chain uses one oracle call and $O(1)$ additional work per transition if a
uniform challenger can be sampled in constant time. The candidate histograms
require $O(Rn)$ storage in a direct implementation. The geometric median in
\eqref{eq:median} is the linear program
\[
 \min_{q\in\Delta_n,z\ge0}\sum_{r,i}z_{ri}
 \quad\text{subject to}\quad
 -z_{ri}\le q_i-\widehat p_r(i)\le z_{ri}.
\]
Thus the statistical upper bound is also polynomial-time. Approximate
minimization to $o(\varepsilon R)$ preserves the guarantee up to a constant
change in accuracy.

\paragraph{Why ordinary confidence is additive.}
For $T$ i.i.d. samples, the empirical TV risk is $O(\sqrt{n/T})$. Changing
one sample changes TV loss by at most $1/T$, so bounded differences adds only
$O(\sqrt{\log(1/\delta)/T})$ to this expectation
\citep{canonne2020short}. This gives the sum
$n+\log(1/\delta)$. Winner-chain states do not have this product
concentration. The half-heavy example in \eqref{eq:half-heavy} retains one
transition for $\Theta(n)$ steps on average. The lower bound below shows that
no different adaptive estimator can bypass the resulting confidence cost.

\section{The Confidence Cost Is Necessary}
\label{sec:lower}

We now show that no adaptive pair selection can attain the additive
confidence dependence of ordinary sampling. For $a\in(0,1)$, define
\begin{equation}
 p_a(1)=a,\qquad p_a(j)=\frac{1-a}{n-1},\quad j\ge2.
 \label{eq:hard-family}
\end{equation}
Take $a_-=1/2-2\varepsilon$ and $a_+=1/2+2\varepsilon$. Then
$\TV(p_{a_-},p_{a_+})=4\varepsilon$.

A query not containing label $1$ returns a fair draw under both hypotheses.
For $\{1,j\}$, let the response be one when the light label wins. Its
Bernoulli parameter is
\begin{equation}
 r(a)=\frac{1-a}{1+(n-2)a}.
 \label{eq:rare-coin}
\end{equation}
On $a\in[1/4,3/4]$, one has $r(a)=\Theta(1/n)$ and
\begin{equation}
 |r'(a)|=\frac{n-1}{(1+(n-2)a)^2}=O(1/n).
 \label{eq:rare-derivative}
\end{equation}
The elementary Bernoulli bound
$D_{\rm KL}(x\|y)\le(x-y)^2/[y(1-y)]$ therefore gives
\begin{equation}
 D_{\rm KL}(\operatorname{Bern}(r(a_-))\| 
 \operatorname{Bern}(r(a_+)))
 \le C_1\frac{\varepsilon^2}{n}.
 \label{eq:one-query-kl}
\end{equation}

For an adaptive learner, condition on each realized transcript before the
next query. Every possible pair obeys \eqref{eq:one-query-kl}, so the KL
chain rule bounds the divergence between length-$Q$ transcript laws by
$C_1Q\varepsilon^2/n$.

If an estimator is $\varepsilon$ accurate with probability $1-\delta$ under
both hypotheses, choosing the closer of $p_{a_-}$ and $p_{a_+}$ produces a
test whose two error probabilities are at most $\delta$. The
Bretagnolle--Huber inequality then requires transcript KL at least
$\log(1/(4\delta))$. Consequently,
\[
 Q\ge c_1\frac{n\log(1/\delta)}{\varepsilon^2},
\]
which proves the lower bound in Theorem~\ref{thm:minimax}.

\section{Bounded Menus Interpolate the Obstruction}
\label{sec:bounded-menu}

The confidence phenomenon is not unique to pairs. It weakens continuously
when a query may expose more light labels at once. Let
$q_{\le k}(n,\varepsilon,\delta)$ denote minimax labeled TV complexity when
the learner may condition on any nonempty menu of size at most $k$.

\begin{proposition}[Bounded-menu confidence lower bound]
\label{prop:bounded-menu}
For $2\le k\le n$ and the parameter ranges of
Theorem~\ref{thm:minimax},
\begin{equation}
 q_{\le k}(n,\varepsilon,\delta)
 =\Omega\!\left(
 \frac{n+(n/k)\log(1/\delta)}{\varepsilon^2}
 \right).
 \label{eq:bounded-menu}
\end{equation}
\end{proposition}

\begin{proof}
The $\Omega(n/\varepsilon^2)$ term already holds with unrestricted
conditional access, since learning all labels is no easier than the standard
multinomial packing used for ordinary samples
\citep{canonne2020short,kleinberg2026providers}.

For confidence, use the family \eqref{eq:hard-family}. A menu excluding label
$1$ is uniform under both hypotheses. Suppose a menu contains label $1$ and
$s\le k-1$ light labels. Conditional on a light outcome, its identity is
uniform under both hypotheses, so all information is in the Bernoulli light
probability
\begin{equation}
 r_s(a)=\frac{s(1-a)}{(n-1)a+s(1-a)}.
 \label{eq:rs}
\end{equation}
Uniformly for $a\in[1/4,3/4]$, one has
$r_s(a)=\Theta(s/n)$ and
\[
 |r_s'(a)|
 =\frac{s(n-1)}{((n-1)a+s(1-a))^2}
 =O(s/n).
\]
The Bernoulli KL is therefore $O(s\varepsilon^2/n)$, and hence at most
$O(k\varepsilon^2/n)$, for every conditional menu and every transcript.
The KL chain rule and Bretagnolle--Huber argument from
Section~\ref{sec:lower} require
$Q=\Omega((n/k)\log(1/\delta)/\varepsilon^2)$. Taking the maximum of the two
lower bounds proves \eqref{eq:bounded-menu}.
\end{proof}

For $k=2$, Theorem~\ref{thm:minimax} matches
\eqref{eq:bounded-menu}. For $k=n$, the lower bound becomes the ordinary
additive rate. This identifies a candidate interpolation law, but we do not
claim its upper bound for intermediate $k$. A direct $k$-winner chain has
stationary law $p$ and a stronger minorization, yet its coordinate Green
function is no longer the minimum-mass graph in Lemma~\ref{lem:profile}.
Determining whether its occupation profile attains \eqref{eq:bounded-menu}
is a separate open problem.

\section{Related Work and Scope}
\label{sec:related}

\paragraph{Restricted conditional learning.}
Conditional sampling was introduced for distribution testing
\citep{canonne2015testing,chakraborty2016power}. The PCOND restriction is
powerful for many label-invariant testing problems, but our target is the
entire labeled distribution. \citet{kleinberg2026providers} characterize
which fixed menu families permit pointwise and uniform labeled learning.
Their global and local winner chains motivate ours. Their global learner
restarts after every terminal sample, while their hierarchical learner keeps
local trajectories to estimate tree splits. Our result instead controls the
full empirical law of one global pairwise trajectory through its complete
mass profile. It both removes the dimension logarithms and supplies the
matching confidence lower bound.

\paragraph{Comparison estimation and sampling.}
Classical BTL work estimates rankings, log qualities, or normalized weights
under passive comparison graphs. Sharp bounds depend on graph Laplacians,
effective resistance, and often bounded dynamic range
\citep{shah2015estimation,negahban2017rank,
agarwal2018accelerated,hendrickx2019resistance,hendrickx2020minimax}.
Our active estimator permits all simplex boundary points and is analyzed in
labeled TV. \citet{fotakis2022perfect} construct exact samples from
exogenously drawn pair comparisons using coupling from the past. Their
distribution-free method first learns every weight to relative accuracy and
solves a different exact-sampling objective.

\paragraph{Biased sampling and Markov concentration.}
Conditional menus are a finite selection-bias model. Classical work proves
identifiability and asymptotic efficiency for fixed biased-sampling designs
\citep{vardi1985empirical,gill1988large}. Generic reversible-chain
concentration scales with a worst-case spectral gap
\citep{jiang2018bernstein}. That route loses a factor $n$ here. The capped
mass calculation is coordinate-specific and nonasymptotic.

\paragraph{Limits and open directions.}
Theorem~\ref{thm:minimax} concerns the family of all pairs. It does not close
the polylogarithmic gaps for every hierarchically comparable menu family.
Proposition~\ref{prop:bounded-menu} gives a broader necessary interpolation,
but whether it is sufficient for every menu cap remains open. The present
algorithm also assumes that every pair is available. Graph-restricted active
designs may require additional connectivity and dynamic-range parameters.

\section{Conclusion}

Pairwise conditional samples retain enough aggregate information to learn a
labeled distribution at the ordinary $n/\varepsilon^2$ rate at constant
confidence. They do not retain the ordinary model's cheap confidence
amplification. Reusing a single winner-chain trajectory attains the optimal
constant-confidence rate, while independent trajectories and a metric median
match the unavoidable multiplicative confidence cost. The proof explains
why a linear relaxation time does not imply a linear loss in estimation
risk: the slow modes and the TV sensitivity are coupled through the target's
capped mass profile.

\bibliography{references}

\end{document}
